% ── CAPÍTULO 4: DISEÑO ALGORÍTMICO ──────────────────────────
\section{Diseño Algorítmico}
\label{sec:disenio}

\subsection{Visión General}

KGeoMIP y KQNodes comparten la misma estructura de alto nivel: (1)~preparar estructuras de datos, (2)~generar particiones candidatas, (3)~evaluar cada candidata con EMD, y (4)~retornar la de menor pérdida. La diferencia fundamental radica en cómo se generan las candidatas y qué oráculo se usa para evaluarlas.

\subsection{Algoritmo GeoMIP (Base)}

GeoMIP construye la tabla de costos navegando el hipercubo por niveles de distancia Hamming.

\begin{algorithm}[H]
\DontPrintSemicolon
\KwIn{NCubos, estado inicial $s_0$, estado final $s_f = 1 - s_0$}
\KwOut{MIP: bi-partición de mínima pérdida}
\BlankLine
$\text{caminos}[0] \leftarrow \{s_0\}$\;
$\text{tabla}[(s_0,s_0)] \leftarrow \mathbf{0}$\;
\Para{$\text{nivel} \leftarrow 1$ \textbf{hasta} $n$}{
    \Para{$s \in \text{caminos}[\text{nivel}-1]$}{
        \Para{$i \leftarrow 1$ \textbf{hasta} $n$}{
            \Si{$s[i] \neq s_f[i]$}{
                $s' \leftarrow s$ con $s'[i] \leftarrow s_f[i]$\;
                \Si{$s'$ no visitado}{
                    calcular $c(s_0, s')$ según ec.\,(4)\;
                    $\text{caminos}[\text{nivel}] \mathrel{+}= \{s'\}$\;
                }
            }
        }
    }
}
$\text{candidatos} \leftarrow$ identificar\_particiones\_optimas()\;
\Para{$(\text{pres}, \text{fut}) \in \text{candidatos}$}{
    $\text{dist} \leftarrow$ bipartir(fut, pres).distribucion\_marginal()\;
    $\text{emd} \leftarrow \EMD(\text{dist},\, p_\text{orig})$\;
}
\KwRet{$\arg\min_\Pi\;\text{emd}[\Pi]$}\;
\caption{GeoMIP: búsqueda geométrico-topológica del bi-MIP.}
\label{alg:geomip}
\end{algorithm}

\subsection{Algoritmo KGeoMIP}

KGeoMIP reutiliza la tabla de costos de GeoMIP y añade generación y evaluación de $k$-particiones candidatas.

\begin{algorithm}[H]
\DontPrintSemicolon
\KwIn{Sistema $V$, número de particiones $k$, presupuesto $B=10{,}000$}
\KwOut{$\kMIP$}
\BlankLine
\_preparar\_geometria()  \tcp*{tabla de costos Hamming, $O(n \cdot 2^n)$}
$v \leftarrow \{(\text{EFECTO},i)\} \cup \{(\text{ACTUAL},d)\}$\;
$\Stirling{|v|}{k} \leftarrow$ calcular según ec.\,(5)\;
\BlankLine
\eSi{$\Stirling{|v|}{k} \leq B$}{
    $C \leftarrow$ \_candidatas\_exactas($k$)   \tcp*{combin. + Hamming + exhaustivo}
}{
    $C \leftarrow$ \_candidatas\_heuristicas($k$) \tcp*{solo geométricas}
}
$C \leftarrow$ deduplica($C$)\;
\Para{$\Pi \in C$}{
    \_evaluar\_k\_particion($\Pi$)\;
}
\KwRet{$\arg\min_\Pi\;\text{memoria}[\Pi].\text{emd}$}\;
\caption{KGeoMIP: búsqueda de $k$-MIP geométrico.}
\label{alg:kgeomip}
\end{algorithm}

\subsubsection{Evaluación de una k-Partición}

\begin{algorithm}[H]
\DontPrintSemicolon
\KwIn{Partición $\Pi = \{G_1,\ldots,G_k\}$}
\KwOut{EMD registrado en memoria}
\BlankLine
\eSi{$k = 2$}{
    $\text{fut} \leftarrow \{i : (\text{EFECTO},i) \in G_1 \cup G_2\}$\;
    $\text{pres} \leftarrow \{d : (\text{ACTUAL},d) \in G_1 \cup G_2\}$\;
    $\text{dist} \leftarrow$ bipartir(fut, pres).distribucion\_marginal()\;
    $\text{emd} \leftarrow \EMD(\text{dist}, p_\text{orig})$\;
}{
    \Para{$G_i \in \Pi$}{
        $f_i \leftarrow \{i:(\text{EFECTO},i)\in G_i\}$;\quad
        $m_i \leftarrow \{d:(\text{ACTUAL},d)\in G_i\}$\;
        $\text{dists}[i] \leftarrow$ bipartir$(f_i,m_i)$.distribucion\_marginal()\;
    }
    $\bar{p} \leftarrow \frac{1}{k}\sum_i \text{dists}[i]$\;
    $\text{emd} \leftarrow \EMD(\bar{p}, p_\text{orig})$\;
}
registrar $(\Pi, \text{emd})$ en memoria\;
\caption{\_evaluar\_k\_particion: cálculo de EMD para una $k$-partición.}
\label{alg:evaluar}
\end{algorithm}

\subsection{Algoritmo QNodes (Base)}

QNodes precomputa la transformada Zeta sobre $\delta = H - p$ y ejecuta el MAO (\textit{Maximum Adjacency Ordering}).

\begin{algorithm}[H]
\DontPrintSemicolon
\KwIn{Sistema preparado}
\KwOut{MIP}
\BlankLine
\_\_preparar\_oraculo()  \tcp*{transf. Zeta: $O(D \cdot N \cdot 2^D)$}
$\text{vert} \leftarrow \{(\text{ACTUAL},d)\} \cup \{(\text{EFECTO},i)\}$\;
\BlankLine
\Para{$v \in \text{vert}$}{
    registrar\_candidato($v$)  \tcp*{pre-pass: singletons}
}
\Mientras{$|\text{vert}| > 1$}{
    $\omega \leftarrow [\text{vert}[0]]$;\quad $\Delta \leftarrow \text{vert}[1:]$\;
    \Mientras{$|\Delta| > 1$}{
        $\delta^* \leftarrow \arg\min_{\delta\in\Delta}\;\bigl[\EMD(\omega \cup \{\delta\}) - \EMD(\{\delta\})\bigr]$\;
        $\omega \mathrel{+}= \delta^*$;\quad $\Delta \mathrel{-}= \delta^*$\;
    }
    registrar\_candidato($\Delta[0]$)  \tcp*{nodo colgante}
    $\text{vert} \leftarrow \omega[{:}-1] \cup \{\omega[-1] \oplus \Delta[0]\}$  \tcp*{contraer}
}
\KwRet{$\arg\min_\Pi\;\text{memoria}[\Pi].\text{emd}$}\;
\caption{QNodes: MAO con oráculo Zeta.}
\label{alg:qnodes}
\end{algorithm}

\subsection{Algoritmo KQNodes}

\begin{algorithm}[H]
\DontPrintSemicolon
\KwIn{Sistema, $k$, presupuesto $B=10{,}000$}
\KwOut{$\kMIP$}
\BlankLine
\_\_preparar\_oraculo()  \tcp*{transformada Zeta}
$\text{vert} \leftarrow$ presente $+$ futuro;\quad $\Stirling{|\text{vert}|}{k} \leftarrow$ calcular\;
\BlankLine
\eSi{$\Stirling{|\text{vert}|}{k} \leq B$}{
    $\Pi^* \leftarrow$ \_exhaustive\_k\_partition(vert, $k$)\;
}{
    $\Pi^* \leftarrow$ \_greedy\_k\_partition(vert, $k$)\;
}
$\delta^* \leftarrow$ \_compute\_partition\_loss($\Pi^*$)\;
\KwRet{Solution(perdida=$\delta^*$, partición=$\Pi^*$)}\;
\caption{KQNodes: búsqueda de $k$-MIP con oráculo Zeta.}
\label{alg:kqnodes}
\end{algorithm}

\subsubsection{Búsqueda Exhaustiva}

\begin{algorithm}[H]
\DontPrintSemicolon
\KwIn{Vértices, $k$}
\KwOut{$k$-partición óptima}
\BlankLine
$\Pi^* \leftarrow \emptyset$;\quad $\delta^* \leftarrow +\infty$\;
\Para{$\Pi \in$ generar\_particiones(vert, $k$)}{
    $\delta \leftarrow$ \_compute\_partition\_loss($\Pi$)\;
    \Si{$\delta < \delta^*$}{
        $\delta^* \leftarrow \delta$;\quad $\Pi^* \leftarrow \Pi$\;
    }
}
\KwRet{$\Pi^*$}\;
\caption{\_exhaustive\_k\_partition: enumera las $\Stirling{n}{k}$ particiones.}
\label{alg:exhaustive}
\end{algorithm}

\subsubsection{Particionamiento Greedy (modo heurístico)}

\begin{algorithm}[H]
\DontPrintSemicolon
\KwIn{Vértices $V$, $k$}
\KwOut{$k$-partición aproximada}
\BlankLine
$\text{grupos} \leftarrow [\,]$;\quad $V_r \leftarrow V$\;
\Para{$i \leftarrow 1$ \textbf{hasta} $k-1$}{
    $\text{mip} \leftarrow$ algorithm($V_r$)  \tcp*{MAO sobre el resto}
    $G \leftarrow$ aplanar(mip)\;
    grupos $\mathrel{+}= G$;\quad $V_r \leftarrow V_r \setminus G$\;
}
grupos $\mathrel{+}= V_r$  \tcp*{último grupo = todo el resto}
\KwRet{grupos}\;
\caption{\_greedy\_k\_partition: heurística greedy por iteraciones del MAO.}
\label{alg:greedy}
\end{algorithm}

\subsection{Estructuras de Datos}

\subsubsection{N-Cubo}

Cada N-Cubo almacena la TPM de un nodo como hipercubo $n$-dimensional. Los accesos usan notación \textit{little-endian}: el bit menos significativo corresponde al último eje.

\begin{lstlisting}[language=Python, caption={Estructura esencial del N-Cubo.}]
@dataclass(frozen=True)
class NCube:
    indice: int               # indice del nodo
    dims: NDArray[np.int8]    # dimensiones activas
    data: np.ndarray          # hipercubo n-dimensional
    memo: dict = field(default_factory=dict)

    def marginalizar(self, ejes: NDArray) -> "NCube":
        resultado = self.data
        for eje in sorted(ejes, reverse=True):
            resultado = resultado.mean(axis=eje)
        return NCube(indice=self.indice, dims=..., data=resultado)
\end{lstlisting}

\subsubsection{Representación de Vértices y Particiones}

Un vértice es la 2-tupla \texttt{(tiempo, índice)} con \texttt{tiempo} $\in \{0=\text{ACTUAL}, 1=\text{EFECTO}\}$. Una $k$-partición es una lista de $k$ grupos, donde cada grupo es una lista de vértices:

\begin{lstlisting}[language=Python, caption={Ejemplo de 3-partición de 4 nodos.}]
particion = [
    [(1,0), (0,0)],                # Grupo 1: Futuro A + Presente A
    [(1,1), (0,1)],                # Grupo 2: Futuro B + Presente B
    [(1,2), (1,3), (0,2), (0,3)], # Grupo 3: resto
]
\end{lstlisting}
